\section{Scope}

This guide outlines the general principles that should be considered and prepared when experimental data is collected for the purpose of publication. It is intended to be read, understood, and applied by all students and researchers performing experimental measurements in the APMS laboratories, including work related to analog, digital, RF, and AI-based systems.

The main stages covered in this guide are:

\begin{enumerate}
\item \textbf{Preparation - before measurements are performed}
\begin{enumerate}
\item Define the quantity or phenomenon to be measured.
\item Establish the expected results or order of magnitude.
\item Identify the required instruments and equipment.
\item Select and document the measurement method.
\item Define how data will be recorded and documented.
\item Ensure that the required instruments are calibrated or verified against a trusted reference.
\end{enumerate}


\item \textbf{During measurements}
\begin{enumerate}
    \item Record measurement data together with time, operator, instrument identification, and relevant environmental conditions.
    \item Verify instrument settings and monitor measurement stability.
    \item Observe and document unexpected behaviour such as drift, instability, or anomalies.
    \item Repeat measurements as necessary to ensure adequate repeatability.
\end{enumerate}

\item \textbf{After measurements are completed}
\begin{enumerate}
    \item Ensure that all data is stored securely and in an accessible format.
    \item Analyse the data for outliers, anomalies, and unexpected trends.
    \item Perform a general estimation of measurement uncertainty.
    \item Review and update the prepared measurement procedure if errors or improvements are identified.
\end{enumerate}

\end{enumerate}

As mentioned, this document outlines the key areas that must be considered. More detailed, hands-on guidance will be referenced where necessary.
